\documentclass[10pt]{letter}
\usepackage[a4paper, margin=0.85in]{geometry}
\usepackage[T1]{fontenc}
\setlength{\parskip}{5pt}
\signature{Harsh Bordekar \\ harshbordekar@gmail.com}
\begin{document}
\begin{letter}{Helsing GmbH \\ Berlin / Munich, Germany}
\opening{Dear Hiring Team,}

I am writing to apply for your AI Research Engineer, ML Engineering position, focused on extending your PyTorch-based deep-learning frameworks and scaling distributed training infrastructure. I am a simulation and machine-learning research engineer with over seven years building physics-based models for aerospace structures, including three years working hands-on across the applied deep-learning stack, building custom model architectures end to end rather than only calling pretrained ones.

At DLR and TU Braunschweig I run large-scale simulation campaigns on HPC/SLURM clusters, generating and managing the high-fidelity, multiscale datasets that my PhD's probabilistic surrogate framework is trained and evaluated against, where job scheduling, turnaround, and data volume are constant practical constraints rather than abstract concerns. I also built a physics-based CNN surrogate for materials-property prediction, wrapped end to end in an agentic, tool-calling interface backed by a hierarchical retrieval pipeline, taking it from data pipeline through model to a deployed, queryable interface rather than stopping at a notebook result. That gives me direct, hands-on experience with the full lifecycle a deep-learning framework has to support, even though my scale to date has been single-node rather than multi-node.

Beyond model-building, I write performance-sensitive C++ for FEM post-processing, and I verify and technically review externally and student-produced simulation models against defined acceptance criteria, tracing failure modes to subtle numerical, meshing, or data-quality issues, a habit of root-cause debugging I would bring directly to your training pipelines. I have also trained and validated a CT-based defect-detection classifier end to end, including its evaluation against ground truth for engineering sign-off.

I also teach: I instruct 60+ Master's students each year in FEM theory and ABAQUS-based composite modelling, covering continuum damage mechanics, progressive damage analysis, and cohesive zone modelling. That has sharpened my ability to build from complex theoretical concepts and explain them clearly to technical and non-technical audiences alike, something I understand matters to how Helsing's engineering culture works internally.

To be transparent about fit: my hands-on scaling experience is with HPC/SLURM job scheduling for large simulation workloads rather than multi-node distributed model training, and I have not yet worked directly with NCCL/MPI, Kubernetes, or Ray. What I bring is a first-principles habit from research, comfort reading and integrating new methods quickly, and genuine motivation to grow specifically into large-scale training infrastructure. I would welcome the opportunity to build that experience directly within your team.

I am motivated by Helsing's mission and by the scale of the engineering challenge behind it, and I would welcome the chance to bring my simulation and applied deep-learning background into your ML engineering team.

Thank you for your consideration. I look forward to discussing how I can contribute.

\closing{Sincerely,}
\end{letter}
\end{document}
